\section{ТЕСТИРОВАНИЕ И ДЕМОНСТРАЦИЯ РАБОТЫ СИСТЕМЫ}

\subsection{Стратегия тестирования}

Тестирование разработанной информационной системы включает модульные тесты, интеграционные тесты, сквозные (E2E) тесты и нагрузочное тестирование. Тесты разделены по компонентам системы. Стратегия тестирования основана на пирамиде тестирования: большое количество быстрых модульных тестов, среднее количество интеграционных тестов и небольшое количество E2E-тестов.

В таблице~\ref{testsummary} приведено общее количество тестов по компонентам.

\begin{table}
  \caption{Общее количество тестов по компонентам}
  \begin{tabular}{|l|l|c|}
    \hline
    \textbf{Компонент} & \textbf{Инструмент} & \textbf{Количество файлов} \\
    \hline
    Backend (Rust) & sqlx::test + tower & 17 \\
    \hline
    Frontend (TypeScript) & Vitest + Testing Library & 10 \\
    \hline
    Analytics-ML (Python) & pytest & 3 \\
    \hline
    Нагрузочное тестирование & k6 & 3 \\
    \hline
    E2E & Playwright & 1 \\
    \hline
  \end{tabular}
  \label{testsummary}
\end{table}

\subsection{Тестирование серверной части}

Интеграционные тесты серверной части реализованы с использованием макроса \texttt{sqlx::test}, автоматически подготавливающего тестовую базу данных с применением миграций. HTTP-запросы выполняются in-process через \texttt{tower::ServiceExt}, что исключает сетевые задержки и позволяет тестировать обработчики в изолированной среде.

Тестовый набор включает 17 файлов, покрывающих:
\begin{enumerate}
    \item аутентификацию и авторизацию --- вход с валидными и невалидными учётными данными, выход, обновление токена, проверку ролей admin/user, ограничение частоты попыток входа (5 за 60 секунд), обязательную смену пароля;
    \item CRUD-операции для всех доменных областей --- животные, надои, воспроизводство, кормление, контакты, местоположения, быки, переводы, запасы, включая пакетные операции и импорт CSV;
    \item генерацию отчётов более 17 типов и аналитику --- KPI, тренды, прогнозы;
    \item интеграцию с Lely --- конфигурация, статус синхронизации, ручной запуск;
    \item граничные случаи --- некорректные данные (невалидный JSON, отрицательные значения, строки превышающей длины); конфликты (дублирование инвентарных номеров); несуществующие ресурсы (404); несанкционированный доступ (401, 403).
\end{enumerate}

Для каждого теста создаётся изолированная тестовая база данных, которая уничтожается после выполнения. Общий вспомогательный модуль (\texttt{common/mod.rs}) предоставляет фабрики для создания тестового состояния, генерации JWT-токенов и построения HTTP-запросов.

Пример структуры интеграционного теста: подготовка (создание тестового приложения с подключением к тестовой БД), выполнение (отправка HTTP-запроса к тестируемому эндпоинту), проверка (asserting статус-кода, структуры JSON-ответа и наличия ожидаемых данных) и очистка (тестовая БД уничтожается автоматически).

Помимо функциональных тестов, реализованы бенчмарки для оценки производительности: SQL-запросов (\texttt{bench\_sql}), генерации PDF (\texttt{bench\_pdf}), аутентификации (\texttt{bench\_auth}) и rate limiting (\texttt{bench\_rate\_limit}). Бенчмарки используют критерий Criterion и запускаются командой \texttt{cargo bench}.

\subsection{Тестирование клиентского приложения}

Модульные тесты клиентского приложения реализованы с использованием Vitest~\cite{Vitest} и @testing-library/jest-dom. Тестовый набор включает 10 файлов, покрывающих UI-компоненты (DataTable --- сортировка, пагинация, рендеринг данных; Modal --- открытие, закрытие, контент; FormField --- валидация, отображение ошибок; Pagination --- рендеринг номеров страниц) и утилиты (валидаторы обязательных полей, минимальной длины и формата email; форматирование дат и чисел; построение запросов с параметрами фильтрации и пагинации; хелперы графиков для конфигурации Chart.js; CRUD-модальные окна с состоянием создания и редактирования; бизнес-правила валидации).

Сквозное (E2E) тестирование реализовано с использованием Playwright. Тесты проверяют ключевые пользовательские сценарии: аутентификация (вход с валидными и невалидными данными), навигация по основным разделам, CRUD-операции (создание, редактирование, удаление животного), экспорт отчётов в CSV и PDF.

\subsection{Тестирование модуля машинного обучения}

Тесты модуля машинного обучения включают: test\_api.py --- тесты API-эндпоинтов (проверка health-эндпоинта, вызов прогнозов для всех 11 моделей с валидными и невалидными входными данными, запуск обучения); test\_models.py --- тесты моделей (обучение на синтетических данных, проверка формата прогноза на соответствие JSON-схеме, проверка доверительных интервалов q10 < q50 < q90); test\_drift\_and\_auth.py --- тесты мониторинга дрейфа (обнаружение аномального распределения прогнозов) и аутентификации по API-ключу (отклонение запросов без ключа или с невалидным ключом).

\subsection{Нагрузочное тестирование}

Нагрузочное тестирование проведено с использованием k6~\cite{K6}. Реализованы три сценария: дымовой тест (smoke.js) с 5 виртуальными пользователями и 10 итерациями для проверки health-эндпоинта и входа в систему; нагрузочный тест (load.js), моделирующий типичную рабочую нагрузку с постепенным увеличением числа пользователей от 10 до 50; стресс-тест (stress.js) для определения предельной нагрузки с увеличением числа пользователей до 200.

Метрики, собираемые в ходе нагрузочного тестирования: время ответа (p50, p90, p95), количество запросов в секунду (RPS), процент ошибок (HTTP 5xx), время ожидания соединения (connection time).

\subsection{Демонстрация системы}

Для демонстрации работы системы используется скрипт \texttt{dev.sh}, автоматически запускающий PostgreSQL, разворачивающий тестовые данные и запускающий все компоненты системы. Тестовые данные включают 100 животных с историей надоев, воспроизводства и кормления за 365 дней.

\subsubsection{Аутентификация}

При первом запуске система создаёт пользователя \texttt{admin} со случайным паролем из 16 символов, который выводится в стандартный поток ошибок. При входе в систему, если установлен флаг обязательной смены пароля, пользователь перенаправляется на форму смены пароля. Минимальная длина пароля --- 8 символов. Демонстрация интерфейса аутентификации представлена на рисунке~\ref{auth-page}.

\begin{figure}
  \centering
  \includegraphics[width=0.95\textwidth,height=0.85\textheight,keepaspectratio]{auth-page-screenshot.png}
  \caption{Страница аутентификации}
  \label{auth-page}
\end{figure}

\subsubsection{Дашборд}

Главная страница системы --- дашборд, отображающий ключевые показатели стада: количество животных (с разбивкой по полу и статусу), средний суточный надой, статус воспроизводства (количество стельных, в охоте, ожидаемых отёлов), количество активных оповещений. Виджеты настраиваются пользователем через панель настроек: можно выбрать отображаемые KPI и их расположение. Пример дашборда представлен на рисунке~\ref{dashboard}.

\begin{figure}
  \centering
  \includegraphics[width=0.95\textwidth,height=0.85\textheight,keepaspectratio]{dashboard-page-screenshot.png}
  \caption{Дашборд системы}
  \label{dashboard}
\end{figure}

\subsubsection{Управление животными}

Страница списка животных предоставляет таблицу с сортировкой по столбцам (кличка, инвентарный номер, пол, дата рождения, статус), фильтрацией (по полу, местоположению, дате рождения) и пагинацией (по 25, 50 или 100 записей). Поддерживается поиск по кличке и инвентарному номеру с использованием нечёткого поиска pg\_trgm (триграммное сходство).

Для каждого животного доступна детальная карточка с полной информацией, разделённой на вкладки: биография (основные данные, происхождение), надои (графики и таблица), воспроизводство (история событий), кормление (план и факт), ветеринарные записи. Поддерживается экспорт карточки в PDF. Пример списка животных представлен на рисунке~\ref{animals-page}, пример детальной карточки --- на рисунке~\ref{animal-card}.

\begin{figure}
  \centering
  \includegraphics[width=0.95\textwidth,height=0.85\textheight,keepaspectratio]{animals-page-screenshot.png}
  \caption{Список животных}
  \label{animals-page}
\end{figure}

\begin{figure}
  \centering
  \includegraphics[width=0.95\textwidth,height=0.85\textheight,keepaspectratio]{animal-card-screenshot.png}
  \caption{Детальная карточка животного}
  \label{animal-card}
\end{figure}

\subsubsection{Отчёты и аналитика}

Система генерирует более 17 типов отчётов (см. таблицу~\ref{reports}), каждый из которых отображается в виде таблицы и/или графика в интерфейсе. Отчёты поддерживают фильтрацию по дате и группе животных, а также экспорт в CSV и PDF.

Аналитическая панель отображает: ключевые показатели эффективности (KPI) --- средний надой, процент стельных, интервал между отёлами; тренды надоев с прогнозом на основе модели Хольта--Винтерса; прогноз стельных коров; кривые лактации по номерам; кластеризацию стада.

\subsubsection{Интеграция с Lely}

При включённой интеграции с Lely система автоматически синхронизирует данные каждые 5 минут. Страница статуса синхронизации отображает дату последней успешной синхронизации для каждого из 23 типов данных и количество синхронизированных записей. Через интерфейс доступен ручной запуск синхронизации и просмотр логов синхронизации.

\subsubsection{Оповещения}

Модуль оповещений автоматически создаёт оповещения на основе настраиваемых пороговых значений. Проверка выполняется каждые 15 минут фоновым заданием. Отслеживаются следующие категории событий:
\begin{enumerate}
    \item падение надоев более чем на 20~процентов за сутки;
    \item высокий \acrshort{scc} (более 200\,000 клеток/мл);
    \item снижение активности более чем на 30~процентов за сутки;
    \item снижение жвачки более чем на 25~процентов за сутки;
    \item ожидаемый отёл (за 7 дней до планируемой даты).
\end{enumerate}

Пороговые значения настраиваются через страницу настроек без перезапуска сервера. Оповещения отображаются в интерфейсе с категоризацией по степени важности (high, medium, low). Поддерживается подтверждение оповещения (acknowledge) для отслеживания реакции фермера. Оповещения также могут быть отправлены через канал уведомлений (браузерные уведомления или Telegram-бот).

\subsection{Мониторинг работоспособности}

Стек Grafana LGTM обеспечивает комплексное наблюдение за системой. Grafana дашборды отображают частоту запросов, время ответа (p50, p95, p99), использование ресурсов (CPU, память) и ошибки по маршрутам. Loki позволяет искать логи по уровню (info, warn, error) и целевому сервису с использованием языка LogQL. Tempo обеспечивает распределённую трассировку запросов через цепочку Nginx $\to$ Backend $\to$ PostgreSQL/ML, позволяя идентифицировать медленные запросы. Prometheus собирает метрики с backend (\texttt{/metrics}), analytics-ml (\texttt{/metrics}) и Redis с интервалом 15 секунд.

Метрики backend включают: \texttt{http\_requests\_total} (счётчик, разбитый по маршруту, методу и статус-коду), \texttt{http\_request\_duration\_seconds} (гистограмма), \texttt{db\_connections\_active}, \texttt{rate\_limit\_rejections\_total}.
